\section*{Guide to the Supplementary Material}
The supplementary material follows the paper's argument, expanding each connection between geometry, learning, and practical allocation. Table~\ref{tab:argument-guide} provides entry points to the complete results and proofs.
\begin{table}[ht]
\centering\small
\caption{\textbf{Where each part of the argument is developed.}}
\label{tab:argument-guide}
\begin{tabular}{@{}p{.36\linewidth}p{.59\linewidth}@{}}
\toprule
Question & Supporting material \\
\midrule
How is allocation value identified? & Fixed-support training (App.~\ref{sec:exact-range-control}); multi-shape factorial (App.~\ref{sec:exact-range-factorial}); frozen controls (App.~\ref{sec:frozen-fixed-support}). \\
How are the constructions derived? & Full-pair geometry and Cosh (App.~\ref{sec:proofs}); TailSpline (App.~\ref{sec:tailspline-details}); BM (App.~\ref{sec:bm-construction}). \\
What are the complete learning and deployment results? & Continuation and MLA (Apps.~\ref{sec:experiment-details},~\ref{sec:mla-results}); adaptation and readout (App.~\ref{sec:mature-details}); frozen profiles and QA (App.~\ref{sec:adjustment-details}); TailSpline (App.~\ref{sec:tailspline-details}). \\
What controls delimit the conclusions? & Range, crossed-table and slot interventions (App.~\ref{sec:compatibility}); profile diagnostics (App.~\ref{sec:profile-diagnostic-audit}); interval design (App.~\ref{sec:interval-design}); video and architecture (Apps.~\ref{sec:video-dit},~\ref{sec:compression-ablation}). \\
\bottomrule
\end{tabular}
\end{table}
\FloatBarrier
